\section*{237塔联合模型完整拆解：运行与规模}
观测时间：2026年11月10日00:54。全区域可行性尚未判定；本次未重启完整优化。
\begin{center}\small\begin{tabular}{p{33mm}p{130mm}}\toprule 对象 & 实测状态\\\midrule
场景 & Rommerskirchen：237塔、2133拍摄任务，2车4机8电池\\
最新运行 & 329a7f263e3525a1\_6948；已退出\\
总耗时／构建 & 663.111515 s／95.606626 s；总耗时含构建与检查\\
退出／搜索节点 & SCIP: error in LP solver!／2；无237塔可行解\\
出错时进程／主机 & RSS 4.213 GiB；可用RAM 4.557 GiB\\
原始变量类型 & 168793二元，12158整数，443925连续\\
约束类型 & 176614线性行，45592非线性行；另40798索引关系\\
本次只读重建 & 72.704256 s；变量和显式行与原运行完全一致\\
\bottomrule\end{tabular}\end{center}\begin{center}\small\begin{tabular}{p{100mm}p{28mm}p{28mm}}\toprule 实际构建模块 & 新增变量 & 新增显式行\\\midrule
任务选择、前后继和航次 & 98614 & 38214\\
实体无人机与电池链 & 10140 & 7032\\
车辆道路阶段 & 47570 & 6150\\
充电事件与账本 & 17316 & 9044\\
充电温度 & 14977 & 1872\\
自由连续航迹与姿态 & 361092 & 33399\\
车机端点、起降与平台互斥 & 40560 & 51636\\
地面相机连续性 & 17628 & 48516\\
能量变量及航次累加 & 15555 & 22110\\
必要界、接触采样、237塔硬要求 & 1424 & 4233\\
合计 & 624876 & 222206\\
\bottomrule\end{tabular}\end{center}
\clearpage
\section*{决策、约束和实际接入方式}
\begin{center}\small\begin{tabular}{p{23mm}p{61mm}p{73mm}}\toprule 模块 & 自由决策／固定输入 & 目前执行方式与缺口\\\midrule
任务与目标 & 237塔完成量；2133任务启用、时刻、前后继 & 整塔9项同时完成；原目标按未完成塔、完工时刻、能耗依次优化。本轮额外固定237塔全部完成。\\
航次与资源 & 156可选槽；任务组合、起落时刻、机／电池身份 & 由平台占用上界得到每车78槽；没有预制固定航次目录。标签对称未处理。\\
车辆道路 & 2×256模式／时长／7次曲线；全部源道路 & 13936节点、20455有向边、104停靠点。原4096动力学采样点初始显式接入0个。\\
无人机轨迹 & 9000段、7次位置／姿态多项式；转移最多3段 & 曲线可优化；大部分动力学、天气、光学和避障在候选出现后验收、加约束。\\
起降协同 & 连续位置、平台时隙、选择停车阶段 & 同一个SCIP模型；当前初始物理点仅8个起落接触点。\\
充电／电池 & 312事件；时刻、所属电池、充电增量 & 同机／同电池前后继与Wh账本；充电温度接入。旧见证没有启用充电。\\
能耗 & 段Wh、任务累加、航次与物理电池余量 & 初始9000段实际积分链接0条；有全局功率／时间下界。候选须通过独立积分区间检查。\\
气象 & 保留预报的时间／高度插值 & 单格点、确定性；6小时可飞窗口；不是空间精细天气或随机/鲁棒优化。\\
地面相机 & 312段相机连续多项式 & 新增17628变量、48516显式行；在还没有区域车辆／航次方案时已经展开。\\
安全与执行 & 地形、塔柱、导线、车机／机机冲突 & 连续候选检查为数值工程验收；非严格定向舍入证明。全区域执行回放及重规划未接入。\\
\bottomrule\end{tabular}\end{center}\small 延迟约束并非被删除，但初始松弛缺少核心关系；“声明了完整验收”不等于“求解已取得全区域可行方案”。
\clearpage
\section*{已定位的问题与证据等级}
\begin{center}\small\begin{tabular}{p{13mm}p{35mm}p{109mm}}\toprule 编号 & 问题 & 证据与影响\\\midrule
F01 & 验收口径错误 & 历史逐层放行只要求完成量大于0，237→24→4仍继续；当前门已纠正，旧全覆盖结论不可用。\\
F02 & 4塔来源 & 初始化器显式取最近4塔、每车2航次；不代表SCIP证明最多完成4塔。\\
F03 & 当前退出原因 & 全覆盖运行663.112s后LP数值错误、节点2、无可行解；出错时可用内存4.557GiB。不是无解证明。\\
F04 & 分母界传播缺失 & T+1-a隐含下界0.001，表达式区间却触及0；SCIP日志把下界改为1e-9。独立LP证明0.001。\\
F05 & 核心约束延迟与弱松弛 & 40798索引关系在选择整数化后才生效；初始4096车辆采样点和9000段能量积分均未显式接入。小反例证实分数点缺少有效下界。\\
F06 & 道路数值域截断 & 256阶段使两车合计仅44/104停靠点可到达并返回；未证明该截断保持原区域可行域。\\
F07 & 曝光约束作用域不一致 & 数学稿仅曝光1ms要求拖影；代码整段20.011s执行。实际光学评价器反例：曝光0.01210px通过，末端2.26443px被拒。\\
F08 & 偏航角定义不一致 & 数学稿psi=atan2(R21,R11)，代码输入heading构造姿态；25度倾角反例输入0度而实际偏航-5.60053度。\\
F09 & 航次标签对称 & 每车78个同质可选槽，未按起飞时刻规范标号，自动对称检测关闭；4个具名航次仅槽位排列可达36072036种。\\
F10 & 全域初始化与验证缺口 & 旧见证0车辆行驶、0充电；全区域仅拍摄就至少21航次、9181.707Wh、充电5674.507Wh。必要耦合未由该见证覆盖。\\
\bottomrule\end{tabular}\end{center}\small F01门控已经修正；其余为现存问题或验证缺口。F03是已观测退出原因；F04、F05、F09是有证据的风险因素，尚未完成逐项消除后的因果归因。
\clearpage
\section*{数值风险与两个可复现的模型差异}
\begin{center}\small\begin{tabular}{p{40mm}p{123mm}}\toprule 检查项 & 实际数值／结论\\\midrule
LP非零系数最小／最大 & 9.912807e-8／2.701087e8\\
系数跨度 & 2.724846e15；这是最大/最小系数比，不是矩阵条件数\\
宽变量界 & 11417个变量绝对界$\ge$1e9；120个$\ge$1e15；96个$\ge$1e18\\
安全分母 & D=T+1-a；T$\ge$0.001a，0$\le$a$\le$1，所以D$\ge$0.001\\
独立LP核验 & 最小D=0.0010000000000000009；optimal\\
日志实际行为 & 分母表达式下界被SCIP从0改成1e-9；倒数允许量级1e9而非已蕴含的1e3\\
分数选择反例 & 两类别均要求x$\ge$10；选择码0.5、x=0仍未触发索引检查；整数原问题最优值10\\
含义 & 索引整数可行性未被此反例否定；缺的是能在分数松弛阶段起作用的有效约束。\\
\bottomrule\end{tabular}\end{center}\begin{center}\small\begin{tabular}{p{40mm}p{123mm}}\toprule 数学合同对照 & 反例结果\\\midrule
拖影：曝光期与整段 & 焦距80mm；7次云台曲线；幅值0.008rad。曝光[1,1.001]s拖影0.012102px，通过实际光学评价器。\\
曝光外末端 & 拖影2.264426px；其余光学检查通过，只有拖影失败。代码将该点也纳入强制验收。\\
云台运动界 & 最大角速度0.039178rad/s、角加速度0.062651rad/s²；均低于现有$\pi$/6上界。\\
偏航角 & 输入heading=0°、机体倾斜25°；构造姿态的atan2(R21,R11)=-5.600529°。\\
\bottomrule\end{tabular}\end{center}\small 以上为独立模块反例，不是237塔可行解。曝光合同与姿态定义须先统一，不能通过修改评价口径宣称已修好。
\clearpage
\section*{道路范围、能量和逐层验收缺口}
\begin{center}\small\begin{tabular}{p{99mm}p{29mm}p{29mm}}\toprule 道路必要条件 & 车1 & 车2\\\midrule
当前阶段数 & 256 & 256\\
可到达并返回的停靠点 & 42 & 44\\
可用道路模式并集 & 6199 & 6497\\
原始阶段×模式组合 & 8830720 & 8830720\\
拓扑前后向必要条件保留组合 & 356136 & 381574\\
使每个停靠点分别可往返的最低阶段数 & 631 & 627\\
\bottomrule\end{tabular}\end{center}\small 两车停靠点并集44/104；被当前阶段数排除60个。631/627仅是分别访问每个点的拓扑下界，既不是全部点同日可访问证明，也不含车辆转弯、时间、碰撞和充电约束。\par\begin{center}\small\begin{tabular}{p{99mm}p{58mm}}\toprule 全237塔必要量 & 下界\\\midrule
全部拍摄驻留时间 & 42683.636800 s（11.856566机时）\\
单航次可用电量 & 438.4 Wh\\
功率下界／航次时长上界 & 774.398550 W／2038.020345 s\\
只计拍摄的航次数 & 至少21次\\
只计拍摄能耗 & 至少9181.707347 Wh\\
初始8电池可用能量 & 3507.2 Wh\\
必须补充的充电量 & 至少5674.507347 Wh；未计转移、爬升等\\
\bottomrule\end{tabular}\end{center}\begin{center}\small\begin{tabular}{p{36mm}p{35mm}p{86mm}}\toprule 历史逐层候选 & 完成塔／任务 & 不能推出的结论\\\midrule
第1—7层 & 237／2133 & 未加入完整物理约束，不证明区域可行\\
第8—9层 & 24／216 & 覆盖下降；不能继续当全覆盖通过\\
第10—11层 & 4／36 & 车辆行驶0、充电0；不验证区域必需的多航次周转\\
\bottomrule\end{tabular}\end{center}
\clearpage
\section*{算法、数据适用性与修复顺序}
\begin{center}\small\begin{tabular}{p{33mm}p{130mm}}\toprule 未闭合部分 & 当前事实\\\midrule
高精度能耗 & M300参考功率=悬停项+阻力项+正爬升项；没有接入此前M100实测模型。推力/加速度改变不直接进入悬停功率项。\\
实际电池运行 & 548Wh名义容量、20\%余量；温度范围内容量常数、冷却0s、待机功率0；充电采用参考分段速率。\\
真实地理与模板 & 塔位/道路/地形来自公开地理源；9项任务、塔柱、导线垂弧、0.5m目标面和车辆参数含工程模板。不是现场逐塔测绘。\\
超图算法 & 自制不精确因子共识，参考2013年ADMM消息思想；不是原作者算法的忠实复现，也没有全局最优保证。\\
历史超图结果 & 200轮／109.298069s；0个验收解、15785条显式约束违反。使用旧模型哈希，不能和当前SCIP据此判优劣。\\
完整超图参数 & 当前n=624876；显式行222206。完整含潜在关系及连续域的m、N\_I、r仍未导出，不能用旧r=78冒充当前完整超图。\\
执行与重规划 & 重规划模块只校验状态、生成输入及触发条件；当前区域求解器不接执行快照。尚无区域闭环执行结果。\\
\bottomrule\end{tabular}\end{center}\begin{center}\small\begin{tabular}{p{12mm}p{70mm}p{75mm}}\toprule 顺序 & 下一项修复 & 验收对象\\\midrule
1 & 统一曝光作用域与偏航定义 & 上述正反例；原合法见证与边界反例都重验\\
2 & 显式正分母、归一化、可证明紧界 & 真实LP行及系数；不再只调求解器参数\\
3 & 道路前后向域传播、必要时细化阶段 & 保留全部地理对象；量化被数值网格排除的决策\\
4 & 提前加入核心下界、消除航次标签对称 & 同数学问题；不能把物理约束全部等到整数候选才加入\\
5 & 建立含行驶/充电/多航次的区域初值 & 237塔全部任务的独立完整验收；失败要反哺同一个联合问题\\
6 & 再比较求解器与超图并执行回放 & 相同模型、相同评价器、有效解质量及总时间；随后联合重规划\\
\bottomrule\end{tabular}\end{center}\footnotesize 本次费用0，无新区域求解、云变更或额外工作者。云接入缺少cloud\_services；布局检查发现既有外部暂存链接，本次未搬动文件。\par\footnotesize 证据：\texttt{research\_quality/inspection\_model\_forensics.json}、\texttt{inspection\_model\_contract\_probes.json}及原始运行manifest。复现脚本：\texttt{diagnose\_inspection\_model\_forensics.py}、\texttt{probe\_inspection\_model\_contract.py}。\par\footnotesize 实现规则参考：\href{https://www.scipopt.org/doc/html/CONS.php}{SCIP约束处理器：负优先级在整数检查之后执行}；\href{https://pyscipopt.readthedocs.io/en/latest/tutorials/lazycons.html}{PySCIPOpt延迟约束：可行性检查与排除违例候选}。
